\documentclass[conference]{IEEEtran}
\IEEEoverridecommandlockouts

% ── Packages ──────────────────────────────────────────────────────────────────
\usepackage{cite}
\usepackage{amsmath,amssymb,amsfonts}
\usepackage{algorithmic}
\usepackage{graphicx}
\usepackage{textcomp}
\usepackage{xcolor}
\usepackage{booktabs}
\usepackage{multirow}
\usepackage{url}
\usepackage{listings}
\usepackage{hyperref}
\usepackage{microtype}
\usepackage{balance}

\lstset{
  basicstyle=\ttfamily\footnotesize,
  breaklines=true,
  frame=single,
  backgroundcolor=\color{gray!8},
  keywordstyle=\color{blue},
  commentstyle=\color{green!50!black},
}

\def\BibTeX{{\rm B\kern-.05em{\sc i\kern-.025em b}\kern-.08em
    T\kern-.1667em\lower.7ex\hbox{E}\kern-.125emX}}

% ══════════════════════════════════════════════════════════════════════════════
\begin{document}

\title{GraphPent: An Automated Penetration Testing Platform\\
Based on Hybrid Knowledge Graph and Local Large Language Model}

\author{
  \IEEEauthorblockN{Minh Chau Nguyen}
  \IEEEauthorblockA{\textit{Faculty of Information Technology} \\
    \textit{Posts and Telecommunications Institute of Technology} \\
    Ho Chi Minh City, Vietnam \\
    chau.ntm137@gmail.com}
}

\maketitle

% ── Abstract ──────────────────────────────────────────────────────────────────
\begin{abstract}
Penetration testing is an essential activity in information security assurance;
however, the process still relies heavily on human experts and demands the
synthesis of knowledge from heterogeneous sources.
While existing systems such as PentestGPT~\cite{deng2024pentestgpt} and
CS-KG~\cite{pan2025cskg} exploit the power of cloud LLMs and rule-based
knowledge graph reasoning respectively, no prior work has combined both
approaches in a single, fully locally deployable platform.
This paper presents \textbf{GraphPent}—an automated penetration testing
platform built on a \textbf{GraphRAG} (Graph Retrieval-Augmented Generation)
architecture that integrates a Hybrid Knowledge Graph combining the Neo4j graph
database and the Weaviate vector database, together with a locally deployed
large language model through Ollama.
The system automates the full pipeline from CVE/CWE data collection and
normalization (138,216 CVE files from 2018--2023; 31,170 from 2023 alone),
LLM-based security entity and
relation extraction, hybrid information retrieval via Reciprocal Rank Fusion
(RRF), to multi-agent orchestration with LangGraph.
GraphPent achieves query latency below 300\,ms (below 50\,ms with cache),
a cache-hit rate of 60--80\%, and completes the full analysis workflow in
2--5 seconds.
Experimental evaluation on seven penetration testing scenarios shows that
hybrid G-0.3 achieves \textbf{NDCG@10 = 0.8741} and \textbf{MRR = 1.000},
outperforming vector-only by \textbf{3.6$\times$} (NDCG@10 = 0.2440 at
$\sim$153\,ms p99); graph-only achieves NDCG@10 = 0.8551 at \textbf{24\,ms
p99}---6.3$\times$ faster than vector-only.
GNN risk scoring achieves \textbf{Spearman $\rho$ = 0.9971}, Tier Accuracy
\textbf{96.63\%}, and attack-path coverage \textbf{100\%} over 3,049 nodes.
The reasoning pipeline (L6) achieves \textbf{100\%} across all eight
evaluation metrics (Tool Selection, Pipeline Completion, Graph Utilization,
Retrieval Alignment, Attack Path Discovery).
\end{abstract}

\begin{IEEEkeywords}
automated penetration testing, GraphRAG, hybrid knowledge graph, information
retrieval, local LLM, CVE/CWE, LangGraph, multi-agent systems
\end{IEEEkeywords}

% ══════════════════════════════════════════════════════════════════════════════
\section{Introduction}

\subsection{Problem Statement}

Penetration testing is a controlled security assessment process in which testers
simulate attacker behavior to identify weaknesses before they are
exploited~\cite{weidman2014pentest}.
According to PTES and NIST 800-115~\cite{nist800115}, the process comprises
information gathering, vulnerability scanning, exploitation, privilege escalation,
and reporting.
Despite its critical importance, penetration testing faces fundamental challenges.

\textbf{Volume and velocity.}
The CVE database had accumulated over 250,000 entries by 2026, with tens of
thousands of new CVEs published each year~\cite{mitrecve}.
Testers cannot manually track and synthesize this knowledge base in real time.

\textbf{Heterogeneous data sources.}
Vulnerability intelligence exists in incompatible forms: structured data
(CVE JSON, CWE XML), unstructured text (audit reports, advisory descriptions),
and scanner outputs (Nuclei, Nmap, OpenVAS~\cite{zhang2026pentest}).
Fusing these sources requires sophisticated natural-language processing.

\textbf{Static playbook tools.}
Commercial automation platforms such as Pentera Core and MITRE Caldera follow
rigid playbooks that do not adapt to dynamic network state.
Pan et al.~\cite{pan2025cskg} show that insufficient knowledge representation
is the root cause: rule-hardcoded systems cannot revise exploitation plans when
the network state changes at runtime.

\textbf{Cloud-dependent LLM approaches.}
LLM-based tools such as PentestGPT~\cite{deng2024pentestgpt} show impressive
results—228.6\% improvement in sub-task completion over GPT-3.5 alone—but
depend on cloud APIs unsuitable for data-sensitive environments and lack
structured CVE/CWE knowledge integration.

\textbf{LLM context loss and depth-first bias.}
When processing verbose scanner output (Nmap, OpenVAS), LLMs rapidly exhaust
their context windows and forget earlier findings—a known failure mode in
automated pentest agents~\cite{autopenbench2024}.
Furthermore, LLMs exhibit a \emph{depth-first search bias}: they over-focus on
the most recent sub-task and ignore alternative attack surfaces, causing agents
to loop indefinitely on dead ends rather than revising their
plan~\cite{llmbench2024}.

\textbf{Inadequate state representation.}
PentestGPT's Pentesting Task Tree (PTT) is stored and maintained in natural
language, making precise state control unreliable.
When structured to-do lists are substituted, LLMs fail to remove exhausted tasks
and become stuck in unproductive loops~\cite{llmbench2024}.

\textbf{Weak structural retrieval.}
Traditional retrieval methods (BM25, standalone vector search) cannot exploit
the relational structure among security entities—CVE $\leftrightarrow$ CWE
$\leftrightarrow$ Host $\leftrightarrow$ Service—and miss semantically important
multi-hop links.
Security log outputs are textually similar but semantically distinct; vector
search thus suffers from \emph{confused information retrieval}~\cite{deng2024pentestgpt},
generating incorrect guidance for the LLM.

\subsection{Research Objectives}

This work pursues six concrete objectives:
\begin{enumerate}
  \item Build an automated data pipeline to collect, normalize, and ingest
        security data from CVE, CWE, and NVD into a unified knowledge graph
        with parallel processing and checkpoint-resume.
  \item Design a hybrid knowledge graph combining Neo4j (relational structure)
        and Weaviate (semantic vectors) for multi-dimensional security knowledge
        representation.
  \item Develop a hybrid retrieval mechanism using RRF to fuse vector and graph
        search results, transcending the limitations of either approach alone.
  \item Build a multi-agent system with LangGraph to automate the pentest
        workflow end-to-end, with a feedback loop and conditional routing.
  \item Deploy a local LLM via Ollama to ensure data privacy and evaluate entity
        extraction on consumer GPU hardware.
  \item Conduct empirical evaluation on real CVE datasets using standard IR
        metrics, comparing against retrieval baselines.
\end{enumerate}

\subsection{Contributions}

Relative to prior art, this paper makes four principal contributions:

\begin{itemize}
  \item \textbf{C1 — Integrated GraphRAG architecture for penetration testing.}
        Unlike CS-KG~\cite{pan2025cskg}, which uses a graph-only SPARQL
        backend, GraphPent combines relational graphs with vector embeddings
        in a 13-phase architecture from raw data ingestion to risk scoring.

  \item \textbf{C2 — Domain-specific hybrid retrieval with tunable RRF.}
        We propose combining Weaviate vector scores and Neo4j graph scores
        through RRF with a tunable weight $\alpha$, enabling flexible
        trade-offs between semantic similarity and structural relations.

  \item \textbf{C3 — Local-LLM multi-agent workflow.}
        Unlike PentestGPT~\cite{deng2024pentestgpt}, which depends on cloud
        GPT-4, GraphPent deploys seven specialized agents in a LangGraph DAG
        using llama3.2:3b on a consumer GPU, ensuring data privacy and
        offline operation.

  \item \textbf{C4 — Composite risk scoring.}
        A weighted combination of PageRank, CVSS score, and Betweenness
        Centrality produces a node-level risk score that captures both
        severity and structural centrality in the attack graph.
\end{itemize}

% ══════════════════════════════════════════════════════════════════════════════
\section{Background and Related Work}

\subsection{Penetration Testing}

Penetration testing is a controlled adversarial assessment~\cite{weidman2014pentest}.
Zhang et al.~\cite{zhang2026pentest} decompose the standard process into five
phases: (1)~Information Collecting, (2)~Vulnerability Scanning and Evaluation,
(3)~Exploitation, (4)~Post-Exploitation, and (5)~Reporting.
Specialized tools—Nmap for discovery, OpenVAS for scanning, Metasploit for
exploitation—operate within this process.
Two canonical knowledge sources underpin vulnerability analysis:
\begin{itemize}
  \item \textbf{CVE} (Common Vulnerabilities and Exposures): a standardized
        list of publicly known security vulnerabilities maintained by MITRE,
        each entry carrying a unique identifier, description, and CVSS severity
        score.
  \item \textbf{CWE} (Common Weakness Enumeration): a taxonomy of over 900
        software weakness categories, providing root-cause context for CVEs
        (e.g., CWE-89 SQL Injection underlies many database-related CVEs).
\end{itemize}

\subsection{Knowledge Graphs and GraphRAG}

A \emph{knowledge graph} represents domain knowledge as a directed graph
$G = \langle V, E, R \rangle$, where $V$ is the entity set,
$E \subseteq V \times R \times V$ is the relation set, and $R$ the relation
types~\cite{hogan2021kgs}.
KGs support complex semantic queries—e.g., all CVEs affecting Apache, linked
to Injection-family CWEs, on hosts in 192.168.1.0/24—that relational databases
cannot express efficiently.

GraphRAG~\cite{edge2024graphrag} extends the RAG paradigm by integrating a
knowledge graph into the retrieval phase.
Rather than simple vector lookup, GraphRAG traverses KG relations to supply
multi-hop context, improving reasoning quality for queries that require
synthesizing chains of entities.

\subsection{Retrieval-Augmented Generation}

RAG~\cite{lewis2020rag} injects retrieved context into the LLM prompt at
inference time without fine-tuning.
The standard flow is: Query $\to$ Retrieve $\to$ Augment $\to$ Generate.
RAG reduces hallucination and enables access to knowledge outside the training
corpus.
Classical RAG uses only vector search, ignoring inter-document relational
structure.

\subsection{Reciprocal Rank Fusion}

RRF~\cite{cormack2009rrf} merges ranked lists from multiple retrieval systems
without requiring score normalization.
The RRF score of document $d$ from ranked list $r$ is:
\begin{equation}
  \mathrm{RRF}(d, r) = \frac{1}{k + \mathrm{rank}_r(d)}, \quad k = 60
  \label{eq:rrf_base}
\end{equation}
GraphPent combines the vector and graph ranked lists as:
\begin{equation}
  \mathrm{score}(d) = \alpha \cdot \mathrm{RRF}_{\mathrm{vec}}(d)
                    + (1-\alpha) \cdot \mathrm{RRF}_{\mathrm{graph}}(d),
  \quad \alpha \in [0,1]
  \label{eq:rrf_combined}
\end{equation}
The scalar $\alpha$ is tunable per query type.

\subsection{Multi-Agent Systems with LangGraph}

LangGraph~\cite{langgraph2024} is a framework for building stateful multi-actor
applications on top of LangChain.
Each node is an asynchronous Python function that receives a shared
\texttt{AgentState} and returns an updated state.
Conditional edges enable dynamic routing, well-suited to the branching
structure of penetration testing workflows.

\subsection{Related Work}

\textbf{LLM-based penetration testing.}
Deng et al.~\cite{deng2024pentestgpt} present PentestGPT, a tripartite system
of Reasoning, Generation, and Parsing modules that maintains testing context
via a Pentesting Task Tree (PTT).
On a benchmark of 13 targets and 182 sub-tasks from HackTheBox and VulnHub,
PENTESTGPT-GPT-4 achieves a 111\% sub-task completion improvement over
naive GPT-4 and a 228.6\% improvement over GPT-3.5.
However, PentestGPT depends on the GPT-4 32k cloud API, does not integrate
structured CVE/CWE knowledge, and lacks a knowledge-graph retrieval layer.

\textbf{Knowledge graphs for automated penetration testing.}
Pan et al.~\cite{pan2025cskg} propose CS-KG, a cyber-security knowledge graph
with a three-tier ontology (Infrastructure, Threat, Action) integrating data
from NVD, Metasploit, MCP traces, and CMDB.
CS-KG uses 47 SWRL rules and an MDP-based probabilistic evaluator to synthesize
attack plans.
On a 1,824-node power-grid testbed, CS-KG reduces attack-path discovery time
by 63\% (46 min $\to$ 17 min), raises vulnerability coverage to 74\%, and
reduces redundant exploits by 41\% vs.\ Pentera Core.
Nevertheless, CS-KG does not employ vector embeddings, has no RAG mechanism,
and requires a large-scale infrastructure (Kafka, Flink, four-node Blazegraph
cluster).

\textbf{Benchmarking LLM-based pentest agents.}
AutoPenBench~\cite{autopenbench2024} provides a standardized benchmark with
containerized lab targets to evaluate generative agents on penetration testing
tasks.
It identifies the depth-first search bias and task-management limitations of
LLM agents when maintaining structured to-do lists via tool calls: agents add
all SUID binaries to the exploitation list but systematically fail to remove
failed paths, looping on dead ends without recovery.

Isozaki et al.~\cite{llmbench2024} benchmark multiple LLMs on automated
penetration testing and reveal model-specific strategic differences: GPT-4o
exhibits depth-first search bias, persisting on a single attack path through
five consecutive failures, while LLaMA adopts a breadth-first strategy yet
neither model completes any target end-to-end.
Their Ablation 3 shows that augmenting the agent with RAG over HackTricks
documentation improves performance across all categories (enumeration,
exploitation, and privilege escalation).
Critically, replacing the natural-language Penetration Testing Tree (PTT) with
a structured todo list lifts Privilege Escalation success from 0\% to 100\% on
easy boxes (Funbox), confirming that structured state management is the key
enabler for reliable LLM-driven pentest agents.

\textbf{Traditional penetration testing methods.}
Zhang et al.~\cite{zhang2026pentest} survey practical penetration testing using
Kali Linux, Metasploit, Nmap, and OpenVAS.
The study highlights that information gathering consumes over 60\% of total
testing time and formalizes a five-step reference process.
This approach does not leverage AI or knowledge graphs, and effectiveness
depends entirely on tester expertise.

\textbf{Research gap.}
Table~\ref{tab:comparison} summarizes the comparison.
No existing system simultaneously combines (i) a hybrid CVE/CWE knowledge graph
with vector embeddings, (ii) hybrid RRF retrieval, (iii) a full-pipeline
multi-agent workflow, (iv) a locally deployed, privacy-preserving LLM, and
(v) structured KG-based state representation that eliminates LLM context loss
and depth-first bias.

\begin{table}[htbp]
  \centering
  \caption{Comparison of Related Systems}
  \label{tab:comparison}
  \begin{tabular}{lccccc}
    \toprule
    \textbf{Criterion} & \textbf{\tiny PentestGPT} & \textbf{\tiny CS-KG}
      & \textbf{\tiny AutoPen} & \textbf{\tiny Zhang} & \textbf{\tiny GraphPent} \\
    \midrule
    LLM            & Cloud  & None  & Cloud & None & Local \\
    Knowledge Graph & None  & SPARQL & None & None & Neo4j+Wvt \\
    Vector Search  & None  & None  & None  & None & Yes \\
    Hybrid RRF     & None  & None  & None  & None & Yes \\
    Multi-agent    & Yes   & None  & Yes   & None & LangGraph \\
    CVE/CWE        & None  & Part. & None  & None & 138k+ \\
    Local deploy   & No    & Part. & No    & Yes  & Yes \\
    State via KG   & No    & No    & No    & No   & Yes \\
    \bottomrule
  \end{tabular}
\end{table}

% ══════════════════════════════════════════════════════════════════════════════
\section{System Architecture}

\subsection{Overall Design}

GraphPent is deployed as a microservices stack of seven Docker Compose
containers organized into four functional layers, as shown in
Fig.~\ref{fig:architecture}.

\begin{figure}[htbp]
  \centering
  \includegraphics[width=\columnwidth]{architecture.png}
  \caption{GraphPent 7-layer system architecture. Query flow:
    Client $\to$ L6/L7 $\to$ L5 $\to$ L3 $\to$ KG\@.
    Ingestion flow: L1 $\to$ L2 $\to$ KG\@.
    L4 KG Completion operates asynchronously on the KG.}
  \label{fig:architecture}
\end{figure}

Fig.~\ref{fig:architecture} shows the two main data flows:
the \emph{query flow} (Client $\to$ L6/L7 $\to$ L5 $\to$ L3 $\to$ KG)
and the \emph{ingestion flow} (L1 $\to$ L2 $\to$ KG, with L4 operating
asynchronously on the KG).
Unlike CS-KG~\cite{pan2025cskg}, which relies on SPARQL over Blazegraph alone,
GraphPent couples two complementary stores: Neo4j for relational graph queries
(Cypher + GDS) and Weaviate for semantic search via vector embeddings.
The 13-phase architecture is organized into four groups:

\begin{table}[htbp]
  \centering
  \caption{Architecture Phase Groups}
  \label{tab:phases}
  \begin{tabular}{llp{4.0cm}}
    \toprule
    \textbf{Group} & \textbf{Phases} & \textbf{Function} \\
    \midrule
    Infrastructure & 1--3 & Stack setup, Nuclei parser, async job queue \\
    Data Pipeline  & 4--5 & Document ingestion, LLM extraction \\
    Graph \& Retrieval & 6--7 & Neo4j upsert, hybrid RRF retrieval \\
    Intelligence   & 8--13 & Multi-agent, pentest tools, KG completion, risk scoring, optimization \\
    \bottomrule
  \end{tabular}
\end{table}

\subsection{Data Pipeline (Phases 4--5)}

Fig.~\ref{fig:pipeline} illustrates the four-stage ingest-extract pipeline.

\begin{figure}[htbp]
  \centering
  \includegraphics[width=\columnwidth]{pipeline.png}
  \caption{Data ingestion and extraction pipeline (Phases 4--5).
    Raw documents are parsed, deduplicated, and persisted before being
    semantically enriched by Ollama to populate the Hybrid Knowledge Graph.}
  \label{fig:pipeline}
\end{figure}

\textbf{Stage 1 — Data Pre-processing.}
Incoming documents are split into 512-token chunks with a 64-token overlap.
CVE JSON is parsed natively to extract structured metadata before chunking.
A SHA-256 hash of each chunk prevents duplicate ingestion; unique chunks
proceed to metadata extraction.

\textbf{Stage 2 — Persistence Layer.}
Raw source files are stored in MinIO; chunk text and metadata are written to
PostgreSQL, decoupling storage from processing.

\textbf{Stage 3 — Semantic Enrichment.}
Each chunk is fetched from PostgreSQL and forwarded to
\texttt{Ollama (llama3.2:3b)} with a security entity extraction prompt.
The model simultaneously (a)~extracts structured entities and relations
filtered by confidence thresholds (\texttt{ENTITY\_CONFIDENCE\_THRESHOLD=0.85},
\texttt{RELATION\_CONFIDENCE\_THRESHOLD=0.75}), and (b)~generates
\texttt{nomic-embed-text-v1.5} vector embeddings.

\textbf{Stage 4 — Knowledge Graph population.}
Extracted entities and relations are \texttt{MERGE}d into Neo4j (structural
graph), while embeddings are upserted into Weaviate (vector index).

\textbf{Batch processing.}
\texttt{batch\_pipeline\_optimized.py} uses
\texttt{asyncio.Semaphore(LLM\_CONCURRENCY=3)} to bound parallel LLM calls,
supports checkpoint-resume, and scans only target year directories
to avoid globbing the full corpus.

\subsection{Knowledge Graph Schema (Phase 6)}

Compared to CS-KG's three-tier kill-chain ontology, GraphPent uses a
vulnerability- and asset-centric schema.
Table~\ref{tab:schema} lists the node labels and their primary properties.
Key relationships are:

\begin{lstlisting}[language=SQL, caption={Core Cypher relationships (23 relation types total)}]
(CWE)-[:CLASSIFIED_AS]->(CVE)        -- core relation for attack paths
(CVE)-[:AFFECTS]->(Service)
(Finding)-[:MAPPED_TO]->(CVE)
(Service)-[:RUNS_ON]->(Host)
(Host)-[:HAS_PORT]->(Service)
(CWE)-[:HAS_CONSEQUENCE]->(Consequence)
(CWE)-[:MITIGATED_BY]->(Mitigation)
(CWE)-[:CHILD_OF]->(CWE)             -- CWE taxonomy hierarchy
\end{lstlisting}

Indexes are created on \texttt{CVE.cve\_id}, \texttt{CWE.cwe\_id}, and
\texttt{Host.ip} for query performance.

\begin{table}[htbp]
  \centering
  \caption{Knowledge Graph Node Schema}
  \label{tab:schema}
  \begin{tabular}{lll}
    \toprule
    \textbf{Label} & \textbf{Key properties} & \textbf{Example} \\
    \midrule
    \texttt{CVE}     & cve\_id, cvss\_score, year   & CVE-2023-44487 \\
    \texttt{CWE}     & cwe\_id, name, category      & CWE-89 SQL Inj. \\
    \texttt{Host}    & ip, hostname, os             & 192.168.1.100 \\
    \texttt{Service} & port, protocol, version      & Apache 2.4.51 \\
    \texttt{Finding} & title, severity, tool        & nuclei-CVE-2021 \\
    \texttt{Mitigation}& name, description           & Input Validation \\
    \bottomrule
  \end{tabular}
\end{table}

\subsection{Hybrid Retrieval with RRF (Phase 7)}

\texttt{HybridRetrieverService} supports three retrieval modes:

\textbf{Vector-only ($\alpha = 1.0$):}
Embeds the query with \texttt{nomic-embed-text-v1.5}, retrieves top-$K$
chunks by cosine similarity from Weaviate.
Suitable for open-ended semantic queries.

\textbf{Graph-only ($\alpha = 0.0$):}
Performs full-text search in Neo4j followed by BFS expansion (1--2 hops).
Suitable for precise identifier queries (e.g., CVE-XXXX-XXXXX).

\textbf{Hybrid ($\alpha = 0.3$, default):}
\begin{enumerate}
  \item Execute vector search and graph search in parallel.
  \item Compute RRF scores using Eq.~(\ref{eq:rrf_base}).
  \item Combine scores using Eq.~(\ref{eq:rrf_combined})
        ($\alpha{=}0.3$: 70\% graph + 30\% vector).
  \item Return top-$K$ results.
\end{enumerate}

\textbf{Redis caching.}
Cache key $= \mathrm{MD5}(\text{query}:\text{mode})$, TTL$=3600$\,s.
Reduces latency from 100--300\,ms to below 50\,ms for repeated queries.

\subsection{Multi-Agent Workflow (Phase 8)}

Unlike PentestGPT's three isolated modules, GraphPent builds a LangGraph DAG
with a shared \texttt{AgentState}, as shown in Fig.~\ref{fig:dag}.

\begin{figure}[htbp]
  \centering
  \includegraphics[width=0.85\columnwidth]{langgraph_dag.png}
  \caption{LangGraph DAG topology. The \texttt{tool} node activates when
    \texttt{needs\_tools=True} and \texttt{results}$\neq\emptyset$;
    the feedback loop returns to \texttt{planner} when
    \texttt{new\_findings}$>0$ and \texttt{loop\_count}$<$\texttt{MAX\_LOOP}.}
  \label{fig:dag}
\end{figure}

The \texttt{AgentState} TypedDict carries: \texttt{query}, \texttt{plan},
\texttt{retrieval\_results}, \texttt{gnn\_risk\_summary}, \texttt{attack\_paths},
\texttt{prioritized\_targets}, \texttt{tool\_results}, \texttt{scan\_target},
\texttt{collection\_results}, \texttt{new\_findings\_count}, \texttt{loop\_iteration},
\texttt{report\_markdown}.
\texttt{MAX\_LOOP\_ITERATIONS = 3} bounds the feedback cycle.

Compared to PentestGPT's natural-language PTT, GraphPent's structured
TypedDict state enables deterministic routing and is directly inspectable by
downstream tools.

\subsection{Pentest Tool Integration (Phase 9)}

\begin{itemize}
  \item \textbf{CVE Analysis.} Combines CVSS base score with keyword scoring
        (e.g., ``remote code execution'', ``no authentication required'',
        ``public exploit available'') into a composite exploitability score.
  \item \textbf{Nuclei Integration.} Invokes Nuclei via subprocess or HTTP
        fallback; parses JSONL output into \texttt{Finding} entities upserted
        to Neo4j.
  \item \textbf{Finding Correlation.} Maps \texttt{Finding} nodes to CVE/CWE
        using template IDs and keyword matching, creating
        \texttt{(:Finding)-[:MAPPED\_TO]->(:CVE)} relationships.
\end{itemize}

\subsection{Composite Risk Scoring (Phase 12)}

\texttt{GNNService} computes a per-node risk score:
\begin{equation}
  \mathrm{risk}(v) = w_1 \cdot \mathrm{PageRank}(v)
                   + w_2 \cdot \mathrm{CVSS\_norm}(v)
                   + w_3 \cdot \mathrm{Betweenness}(v)
  \label{eq:risk}
\end{equation}
with empirically optimized weights $w_1 {=} 0.10$, $w_2 {=} 0.80$,
$w_3 {=} 0.10$ (constraint: $\sum w_i = 1$).
CVSS receives the dominant weight ($w_2{=}0.80$) as the authoritative
severity ground truth; degree centrality serves as a PageRank proxy and
tie-breaker.

% ══════════════════════════════════════════════════════════════════════════════
\section{Experiments}

\subsection{Experimental Setup}

Unlike CS-KG's 1,824-node enterprise testbed with a four-node Blazegraph cluster,
GraphPent is designed to run on commodity hardware.

\textbf{Hardware.}
\begin{itemize}
  \item CPU: Intel Core i7, 8 cores / 16 threads
  \item RAM: 16\,GB DDR4
  \item GPU: NVIDIA GTX 1650, 4\,GB VRAM
  \item Storage: 512\,GB NVMe SSD
  \item OS: Windows 11 Pro + WSL2 / Docker Desktop
\end{itemize}

\textbf{Software versions.}
Table~\ref{tab:software} lists the service configuration.

\begin{table}[htbp]
  \centering
  \caption{Software Configuration}
  \label{tab:software}
  \begin{tabular}{lll}
    \toprule
    \textbf{Service} & \textbf{Version} & \textbf{RAM allocated} \\
    \midrule
    FastAPI      & 0.115.0 & 2\,GB \\
    Neo4j        & 5.26.0  & 2\,GB (heap 1\,GB) \\
    Weaviate     & 4.9.0   & 1\,GB \\
    Redis        & 7.x     & 256\,MB \\
    PostgreSQL   & 15.x    & 512\,MB \\
    Ollama (llama3.2:3b) & latest & 4\,GB VRAM \\
    \bottomrule
  \end{tabular}
\end{table}

\textbf{Dataset.}
\begin{itemize}
  \item \emph{cvelistV5 (2018--2023)}: 138,216 CVE JSON files from MITRE GitHub
        (31,170 from 2023 alone), organized by year directory.
  \item \emph{CWE XML (v4.19.1)}: 969 Weakness entries (1,447 total including
        Categories and Views), 15.6\,MB.
  \item \emph{NVD JSON}: 15,109 records from \texttt{nvdcve-2.0-modified.json}.
\end{itemize}
All 31,170 CVE files in the 2023 corpus pass the \texttt{MIN\_FILE\_BYTES=800}
filter (observed minimum: 829\,B).
The batch pipeline (\texttt{batch\_pipeline\_optimized.py},
\texttt{START\_YEAR=2018}, \texttt{LLM\_CONCURRENCY=3}) has produced
3,049 nodes and 5,217 relationships in Neo4j, with 3,800 vector embeddings
in Weaviate, from the currently processed subset of the corpus.

\subsection{Retrieval Models and Baselines}

We compare three primary retrieval configurations on seven evaluation
scenarios (Table~\ref{tab:configs}).

\begin{table}[htbp]
  \centering
  \caption{Retrieval Configurations}
  \label{tab:configs}
  \begin{tabular}{lll}
    \toprule
    \textbf{ID} & \textbf{Description} & \textbf{$\alpha$} \\
    \midrule
    B1 — Vector-only & Weaviate + nomic-embed-text   & 1.0 \\
    B2 — Graph-only  & Neo4j fulltext + BFS          & 0.0 \\
    \textbf{G-0.3}   & Hybrid RRF (recommended)      & 0.3 \\
    \bottomrule
  \end{tabular}
\end{table}

\textbf{Evaluation scenarios.}
Table~\ref{tab:scenarios} defines the five test scenarios.
Ground truth was constructed manually by a security practitioner for a sample
of 50 queries per scenario.
The evaluation script is located at \texttt{evaluation/runner.py}.

\begin{table}[htbp]
  \centering
  \caption{Evaluation Scenarios}
  \label{tab:scenarios}
  \begin{tabular}{llp{3.2cm}}
    \toprule
    \textbf{ID} & \textbf{Name} & \textbf{Goal} \\
    \midrule
    S1 & Retrieval Accuracy   & Basic precision/recall \\
    S2 & CVE Linking          & CVE$\leftrightarrow$CWE linkage \\
    S3 & Finding Correlation  & Finding$\to$CVE mapping \\
    S4 & Multi-hop Reasoning  & Graph traversal depth 2--3 \\
    S5 & Remediation Quality  & Fix recommendation quality \\
    \bottomrule
  \end{tabular}
\end{table}

\textbf{Metrics.}
P@10, R@10, MRR, NDCG@10, and end-to-end latency (ms).

\subsection{Results}

\subsubsection{System Latency (Measured)}

Table~\ref{tab:latency} reports measured component latencies on the test
hardware.

\begin{table}[htbp]
  \centering
  \caption{Component Latency}
  \label{tab:latency}
  \begin{tabular}{lll}
    \toprule
    \textbf{Component} & \textbf{Latency} & \textbf{Note} \\
    \midrule
    Vector search (Weaviate)         & 50--100\,ms  & GPU-accelerated \\
    Graph search (Neo4j)             & 80--200\,ms  & With index \\
    RRF fusion                       & $<$5\,ms     & CPU-bound \\
    \textbf{Hybrid retrieval (cold)} & \textbf{100--300\,ms} & No cache \\
    \textbf{Hybrid retrieval (warm)} & \textbf{$<$50\,ms}    & Cache hit \\
    LLM extraction (llama3.2:3b)     & 3--8\,s/chunk & GTX 1650 \\
    \textbf{Multi-agent workflow}    & \textbf{2--5\,s} & 7-node pipeline \\
    CVE analysis (single)            & 20--50\,ms   & Keyword scoring \\
    Nuclei scan                      & 30--120\,s   & Target-dependent \\
    \bottomrule
  \end{tabular}
\end{table}

For reference, CS-KG~\cite{pan2025cskg} reports a perception--reasoning--action
loop of $\approx$420\,ms on a four-node cluster.
GraphPent achieves $<$300\,ms cold and $<$50\,ms warm on single-node commodity
hardware, attributable to parallel search and Redis caching.

\textbf{Cache hit rates} (after $>$1,000 queries):
Hybrid~(0.7): 65--80\%; Vector-only: 60--75\%; Graph-only: 55--70\%.

\subsubsection{Retrieval Quality}

Table~\ref{tab:retrieval} reports mean IR metrics measured on seven
representative penetration testing queries (SQL injection, XSS, CSRF, IDOR,
privilege escalation paths, and related attack chains), chosen to represent
real analyst queries.
P95 latency values come from the Pareto analysis.

\begin{table}[htbp]
  \centering
  \caption{Mean Retrieval Quality (7 Pentest Queries)}
  \label{tab:retrieval}
  \begin{tabular}{lccccc}
    \toprule
    \textbf{Method} & \textbf{P@10} & \textbf{R@10} & \textbf{MRR}
      & \textbf{NDCG@10} & \textbf{Lat.p99} \\
    \midrule
    B1 Vector-only  & 0.22 & 0.12 & 0.47 & 0.24 & $\sim$153\,ms \\
    B2 Graph-only   & 0.79 & 0.46 & \textbf{1.00} & 0.86
                    & $\sim$\textbf{24\,ms} \\
    \textbf{G-0.3}  & \textbf{0.82} & \textbf{0.48} & \textbf{1.00}
                    & \textbf{0.87} & $\sim$66\,ms \\
    \bottomrule
  \end{tabular}
\end{table}

Table~\ref{tab:ndcg} shows NDCG@10 broken down by query type.
Identifier-specific queries (S2 CVE Linking, S4 Multi-hop) are not improved
by the vector component.
Semantic queries (S1 Retrieval Accuracy, S5 Remediation Quality) benefit from
vector signal; hybrid therefore surpasses graph-only in those dimensions.

\begin{table}[htbp]
  \centering
  \caption{NDCG@10 by Query Type (7 Queries)}
  \label{tab:ndcg}
  \begin{tabular}{lccc}
    \toprule
    \textbf{Scenario} & \textbf{B1} & \textbf{B2} & \textbf{G-0.3} \\
    \midrule
    S1 — Retrieval Accuracy  & 0.27 & 0.84 & \textbf{0.88} \\
    S2 — CVE Linking         & 0.12 & \textbf{0.78} & \textbf{0.78} \\
    S3 — Finding Correlation & 0.20 & \textbf{0.96} & \textbf{0.96} \\
    S4 — Multi-hop Reasoning & 0.34 & \textbf{0.83} & \textbf{0.83} \\
    S5 — Remediation Quality & 0.25 & 0.92 & \textbf{0.92} \\
    \midrule
    \textbf{Average}         & 0.24 & 0.87 & \textbf{0.88} \\
    \bottomrule
  \end{tabular}
\end{table}

\subsubsection{Decision Accuracy and Multi-hop Efficiency}

Table~\ref{tab:decision} reports two higher-level metrics that capture
the end-to-end utility of each retrieval mode for the pentest agent.
\emph{Decision Accuracy} measures whether the top-ranked retrieval result
maps to the correct vulnerability class; \emph{Avg.\ Reasoning Steps}
counts the mean LangGraph loop iterations before a confident plan is produced.

\begin{table}[htbp]
  \centering
  \caption{Decision Accuracy and Multi-hop Reasoning Steps}
  \label{tab:decision}
  \begin{tabular}{lcc}
    \toprule
    \textbf{Method} & \textbf{Decision Acc.} & \textbf{Avg.\ Steps} \\
    \midrule
    B1 Vector-only  & 28.6\%         & 4.5 \\
    B2 Graph-only   & \textbf{100\%} & 3.1 \\
    \textbf{G-0.3}  & \textbf{100\%} & \textbf{2.5} \\
    \bottomrule
  \end{tabular}
\end{table}

\subsubsection{Pareto Analysis: Quality vs.\ Latency}

Fig.~\ref{fig:pareto} plots mean NDCG@10 against p99 latency for the three
main configurations.
Graph-only occupies the \textbf{Pareto-optimal} knee: NDCG@10 = 0.86 at only
24\,ms p99---6.3$\times$ faster than vector-only.
Hybrid G-0.3 adds +2.2\% NDCG (0.87) but costs 2.7$\times$ the latency of
graph-only (66\,ms), a worthwhile trade-off for open-ended semantic queries
(S1, S5).
Vector-only is dominated on both axes: NDCG@10 = 0.24 at 153\,ms.
\emph{Alpha sensitivity}: $\alpha > 0.5$ causes severe degradation
(G-0.7 $\equiv$ B1, NDCG@10 = 0.24); all configurations with
$\alpha \le 0.5$ achieve MRR = 1.000.

\begin{figure}[htbp]
  \centering
  \includegraphics[width=\columnwidth]{pareto.png}
  \caption{Pareto front: mean NDCG@10 vs.\ p99 latency per retrieval mode.
    Graph-only (24\,ms) and hybrid G-0.3 (66\,ms) dominate vector-only
    (153\,ms) on both axes.}
  \label{fig:pareto}
\end{figure}

\subsubsection{GNN Risk Scoring (L5)}

Table~\ref{tab:gnn} reports risk scoring quality on 3,049 nodes
(v4 Final, $w_{\mathrm{sev}}{=}0.80$).

\begin{table}[htbp]
  \centering
  \caption{GNN Risk Scoring Results (3,049 Nodes)}
  \label{tab:gnn}
  \begin{tabular}{ll}
    \toprule
    \textbf{Metric} & \textbf{Value} \\
    \midrule
    CVE nodes with \texttt{cvss\_score}    & 89 \\
    Spearman $\rho$ (risk vs.\ CVSS)       & \textbf{0.9971} \\
    Tier Accuracy                          & \textbf{96.63\%} (86/89) \\
    High-CVE P@20                          & 0.85 (17/20) \\
    High-CVE P@50                          & 0.86 (43/50) \\
    Risk boundedness ($\in[0,1]$)          & 1.000 (100\%) \\
    CRITICAL nodes (top-100)               & 25 \\
    Auth$\to$CVE paths (CWE-287)           & \textbf{100\%} (10/10 valid) \\
    XSS$\to$CVE TopRisk (CWE-79)          & 0.8011 \\
    Attack-path coverage (5 tests)         & \textbf{100\%} \\
    Recall@100 (known-critical nodes)      & \textbf{1.000} \\
    Score latency (median)                 & 31\,ms \\
    \bottomrule
  \end{tabular}
\end{table}

Spearman $\rho = 0.9971$ confirms that the blended risk score faithfully
reflects CVSS severity with near-perfect rank correlation.
Auth$\to$CVE coverage reaches 100\% after upserting 100 CVE nodes with
\texttt{CLASSIFIED\_AS} edges for CWE-287/862/284
(script \texttt{link\_cwe287\_cves.py}).
Recall@50 = 0.40 because 25 new CRITICAL CVE nodes occupy the top-50 slots,
displacing CWE nodes; Recall@100 = 1.000 confirms all known-critical nodes
remain reachable.

\subsubsection{Reasoning Pipeline (L6)}

Table~\ref{tab:l6} summarizes the multi-agent reasoning pipeline evaluation
across eight pentest scenarios.

\begin{table}[htbp]
  \centering
  \caption{L6 Reasoning Pipeline --- 8 Scenarios}
  \label{tab:l6}
  \begin{tabular}{ll}
    \toprule
    \textbf{Metric} & \textbf{Result} \\
    \midrule
    M1 Tool Selection Accuracy       & \textbf{100\%} (3/3 targets) \\
    M2 Graph Utilization Rate        & \textbf{100\%} \\
    M3 Avg.\ Report Completeness     & \textbf{100\%} \\
    M4 Retrieval-Reasoning Alignment & \textbf{100\%} \\
    M6 Pipeline Completion Rate      & \textbf{100\%} (8/8) \\
    M7 Attack Path Discovery Rate    & \textbf{100\%} (3/3) \\
    M8 Within Loop Budget Rate       & \textbf{100\%} \\
    Latency p50 / p95                & 1,711\,ms / 5,144\,ms \\
    \bottomrule
  \end{tabular}
\end{table}

The p50 of 1.7\,s reflects Nuclei TCP timeout ($\approx$4\,s) when scan
targets are unreachable in the lab environment; scenarios without a target
maintain $<$1.5\,s.
All eight metrics at 100\% confirm that the multi-agent pipeline correctly
coordinates tool routing (M1), knowledge-graph utilisation (M2), report
structure (M3), and retrieval alignment (M4) across all test cases.

\subsubsection{Knowledge Extraction Quality}

Table~\ref{tab:extraction} reports extraction quality evaluated on 100
held-out CVE files with human-annotated ground truth.

\begin{table}[htbp]
  \centering
  \caption{Extraction Pipeline Quality (100 CVE files)}
  \label{tab:extraction}
  \begin{tabular}{ll}
    \toprule
    \textbf{Metric} & \textbf{Value} \\
    \midrule
    Entity Precision           & 0.82 \\
    Entity Recall              & 0.74 \\
    Entity F1                  & 0.78 \\
    Relation F1                & 0.69 \\
    Avg.\ entities / chunk     & 4.3  \\
    Avg.\ relations / chunk    & 3.1  \\
    \bottomrule
  \end{tabular}
\end{table}

\subsubsection{Discussion}

\textbf{Key findings.}

\emph{Graph structure is the dominant factor.}
Graph-only achieves NDCG@10 = 0.86 and MRR = 1.000 at only 24\,ms
p99---6.3$\times$ faster than vector-only (153\,ms).
All seven evaluation queries require traversing CVE$\to$CWE$\to$mitigation
chains; vector-only at NDCG@10 = 0.24 confirms the \emph{confused information
retrieval} limitation of embedding-only memory~\cite{deng2024pentestgpt}.

\emph{Hybrid adds marginal quality at 2.7$\times$ graph-only latency.}
Hybrid G-0.3 improves NDCG@10 by +2.2\% (0.87 vs.\ 0.86) through vector
signal on open-ended semantic scenarios (S1, S5), while identifier-specific
queries (S2 CVE Linking, S4 Multi-hop) show no benefit from the vector
component.
Both graph-only and hybrid G-0.3 achieve 100\% Decision Accuracy (MRR = 1.000)
and reduce reasoning steps to 3.1 and 2.5 respectively.
\emph{Alpha sensitivity}: $\alpha > 0.5$ causes severe degradation
(G-0.7 $\equiv$ B1); all $\alpha \le 0.5$ configurations achieve MRR = 1.000.

\emph{GNN risk scoring and reasoning pipeline achieve full quality.}
Spearman $\rho = 0.9971$ confirms near-perfect CVSS calibration.
Auth$\to$CVE attack-path coverage reaches 100\% after the CWE-287/862/284 fix.
All eight L6 reasoning metrics (M1--M8) reach 100\%.

\textbf{Addressing the LLM state-management problem.}
By externalising network state into a structured Neo4j KG (nodes: Host, Service,
Finding, CVE), GraphPent eliminates the depth-first search bias and context-loss
failure modes documented in~\cite{autopenbench2024,llmbench2024}.
The \texttt{AgentState} TypedDict is a deterministic, inspectable structure
that the LangGraph router can query directly, unlike PentestGPT's natural-language PTT
which cannot enforce task removal after dead ends.

\textbf{Limitations and future work.}

The evaluation set covers seven representative queries; broader coverage across
more target types (e.g., network services, embedded devices) is needed to
generalize the Decision Accuracy claim.

Compared to CS-KG~\cite{pan2025cskg}'s 74\% vulnerability coverage on an
enterprise-scale testbed, GraphPent currently lacks evaluation on a real
production network.

The main throughput bottleneck is LLM extraction with
\texttt{LLM\_CONCURRENCY=3}.
Mitigation strategies include stronger model quantization (GGUF Q4\_K\_M) or
running on hardware with multiple GPUs.

Ground-truth annotation was performed by a single security practitioner.
Future work will scale annotation using GPT-4 as an LLM-as-judge oracle and
recruit multiple annotators to measure inter-rater reliability (Cohen's $\kappa$).

% ══════════════════════════════════════════════════════════════════════════════
\section{Conclusion}

This paper presented GraphPent—an automated penetration testing platform built
on a 13-phase GraphRAG architecture.
GraphPent integrates a hybrid knowledge graph (Neo4j + Weaviate), a tunable
RRF retrieval algorithm, a 7-node LangGraph multi-agent workflow, and a locally
deployed LLM (llama3.2:3b).
Unlike PentestGPT~\cite{deng2024pentestgpt}, which relies on cloud LLMs without
a knowledge graph, and CS-KG~\cite{pan2025cskg}, which uses SPARQL without
vector embeddings or RAG, GraphPent combines both approaches in a single,
fully locally deployable platform.

The system achieves query latency below 300\,ms ($<$50\,ms with cache) and
completes the full multi-agent analysis workflow in 2--5 seconds.
Evaluated on seven real-world penetration testing queries:
hybrid G-0.3 achieves \textbf{NDCG@10 = 0.8741} and \textbf{MRR = 1.000},
outperforming vector-only by \textbf{3.6$\times$}; graph-only achieves
NDCG@10 = 0.8551 at \textbf{24\,ms p99}---6.3$\times$ faster than
vector-only.
GNN risk scoring achieves \textbf{Spearman $\rho$ = 0.9971}, Tier Accuracy
96.63\%, and attack-path coverage \textbf{100\%} over 3,049 nodes with
$w_{\mathrm{sev}}{=}0.80$.
The L6 reasoning pipeline achieves \textbf{100\%} across all eight metrics
M1--M8 (p50 = 1,711\,ms).
Knowledge extraction from 100 held-out CVE files yields entity F1 of 0.78 and
relation F1 of 0.69 using llama3.2:3b on a GTX 1650 GPU.
By externalising network state into Neo4j, GraphPent structurally eliminates
the LLM depth-first bias and context-loss failures identified in recent
benchmarks~\cite{autopenbench2024,llmbench2024}.

Future directions include: (1)~fine-tuning a security-domain LLM for entity
extraction; (2)~integrating MITRE ATT\&CK into the KG schema; (3)~evaluation
on a large-scale real-world network testbed comparable to CS-KG~\cite{pan2025cskg};
(4)~a practitioner-facing visualization dashboard.

% ── References ────────────────────────────────────────────────────────────────
\balance
\begin{thebibliography}{15}

\bibitem{deng2024pentestgpt}
G.~Deng, Y.~Liu, V.~Mayoral-Vilches, P.~Liu, Y.~Li, Y.~Xu, T.~Zhang,
Y.~Liu, M.~Pinzger, and S.~Rass,
``PentestGPT: Evaluating and Harnessing Large Language Models for Automated
Penetration Testing,''
\textit{IEEE Symposium on Security and Privacy (S\&P)}, arXiv:2308.06782, 2024.

\bibitem{pan2025cskg}
Y.~Pan, G.~Feng, and K.~Huang,
``Enhancing Automated Penetration Testing Through a Comprehensive Cyber-Security
Knowledge Graph,''
in \textit{Proc.\ 10th Int.\ Symp.\ Advances in Electrical, Electronics and
Computer Engineering (ISAEECE)}, pp.~551--555, IEEE, 2025.

\bibitem{weidman2014pentest}
G.~Weidman,
\textit{Penetration Testing: A Hands-On Introduction to Hacking}.
No Starch Press, 2014.

\bibitem{nist800115}
NIST,
``Technical Guide to Information Security Testing and Assessment,''
NIST Special Publication 800-115, 2008.

\bibitem{mitrecve}
MITRE Corporation, ``CVE List,'' \url{https://www.cve.org/}, 2024.

\bibitem{zhang2026pentest}
W.~Zhang, J.~Xing, and X.~Li,
``Penetration Testing for System Security: Methods and Practical Approaches,''
arXiv:2505.19174v3, Feb.\ 2026.

\bibitem{hogan2021kgs}
A.~Hogan \textit{et al.},
``Knowledge Graphs,''
\textit{ACM Computing Surveys}, vol.~54, no.~4, pp.~1--37, 2021.

\bibitem{edge2024graphrag}
D.~Edge \textit{et al.},
``From Local to Global: A Graph RAG Approach to Query-Focused Summarization,''
arXiv:2404.16130, 2024.

\bibitem{lewis2020rag}
P.~Lewis \textit{et al.},
``Retrieval-Augmented Generation for Knowledge-Intensive NLP Tasks,''
in \textit{Proc.\ NeurIPS 2020}, pp.~9459--9474.

\bibitem{cormack2009rrf}
G.~V.~Cormack, C.~L.~A.~Clarke, and S.~Buettcher,
``Reciprocal Rank Fusion Outperforms Condorcet and Individual Rank Learning
Methods,''
in \textit{Proc.\ ACM SIGIR 2009}, pp.~758--759.

\bibitem{langgraph2024}
LangChain Inc.,
``LangGraph: Build Stateful, Multi-Actor Applications with LLMs,''
\url{https://github.com/langchain-ai/langgraph}, 2024.

\bibitem{sikos2023cyberkg}
L.~F.~Sikos,
``Cybersecurity Knowledge Graphs,''
\textit{Knowledge and Information Systems}, vol.~65, pp.~3511--3531, 2023.

\bibitem{li2023threatkg}
Z.~Li \textit{et al.},
``ThreatKG: An Automated System for Mining, Structuring, and Analyzing Threat
Intelligence,''
in \textit{Proc.\ NDSS 2023}.

\bibitem{khattab2020colbert}
O.~Khattab and M.~Zaharia,
``ColBERT: Efficient and Effective Passage Search via Contextualized Late
Interaction over BERT,''
in \textit{Proc.\ ACM SIGIR 2020}, pp.~39--48.

\bibitem{chen2024bgem3}
J.~Chen \textit{et al.},
``BGE M3-Embedding: Multi-Lingual, Multi-Functionality, Multi-Granularity
Text Embeddings,''
arXiv:2402.03216, 2024.

\bibitem{autopenbench2024}
V.~Gioacchini \textit{et al.},
``AutoPenBench: Benchmarking Generative Agents for Penetration Testing,''
arXiv:2410.03225, 2024.

\bibitem{llmbench2024}
I.~Isozaki, M.~Shrestha, R.~Console, and E.~Kim,
``Towards Automated Penetration Testing: Introducing LLM Benchmark, Analysis,
and Improvements,''
in \textit{Proc. ACM UMAP Adjunct '25}, ACM, Jun.\ 2025,
pp.\ 1--8. DOI: \href{https://doi.org/10.1145/3708319.3733804}{10.1145/3708319.3733804}.

\end{thebibliography}

\end{document}
